\pagenumbering{arabic}\setcounter{page}{1}
\chapter{PENDAHULUAN}

\section{Latar Belakang}

Beberapa tahun terakhir ditandai oleh perkembangan kecerdasan buatan
(\textit{artificial intelligence}, AI) yang berlangsung sangat pesat, baik
pada layanan komersial berbayar maupun pada model \textit{open-weight}.
Kapasitas komputasi AI global tumbuh rata-rata 3,3 kali per tahun sejak 2022
dan mencapai 17,1 juta \textit{H100-equivalents} \citep{aiindex2026},
sementara biaya inferensi untuk kemampuan yang setara dengan GPT-3.5 turun
lebih dari 280 kali antara November 2022 dan Oktober 2024 \citep{aiindex2025}.

\textit{Large Language Model} (LLM) menghasilkan teks secara autoregresif.
Pada setiap langkah, model menghitung distribusi peluang atas seluruh
kosakata, memilih satu \textit{token} diskret dari distribusi tersebut,
menambahkannya ke rangkaian masukan, kemudian mengulangi proses yang sama
untuk menghasilkan \textit{token} berikutnya \citep{vaswani}. Kemampuan
penalaran model juga dibangun di atas mekanisme tersebut. Teknik
\textit{chain-of-thought} mendorong model untuk menghasilkan langkah-langkah
penalaran dalam bentuk teks sebelum menghasilkan jawaban akhir
\citep{wei}, sehingga setiap langkah penalaran tetap diwujudkan melalui
pemilihan \textit{token} diskret.

Peran LLM kemudian berkembang dari sistem percakapan menjadi \textit{agent}.
Model kini dapat melakukan pencarian informasi, terhubung dengan sumber data
dan lingkungan eksternal melalui \textit{Model Context Protocol} (MCP),
memanggil \textit{tools}, hingga mengoperasikan antarmuka komputer.
\textit{Agent} merupakan sistem yang menerima suatu tujuan, menyusun rencana,
dan melakukan tindakan sesuai lingkungan secara relatif otonom
\citep{russellnorvig}.

Beberapa agen kemudian dapat digabungkan ke dalam satu \textit{Multi-Agent
System} (MAS). Setiap agen dapat diberikan peran yang berbeda dan saling
bertukar informasi untuk menyelesaikan suatu tugas, sehingga keseluruhan
sistem dapat dipandang sebagai suatu bentuk kolaborasi antarsistem yang
menyerupai tim kerja \citep{wooldridge}. MAS berbasis LLM telah digunakan
dalam berbagai tugas, termasuk penalaran matematika dan sains,
\textit{question answering}, penggunaan \textit{tools}, hingga pengambilan
keputusan \citep{guo}.

Peningkatan kemampuan tersebut juga diikuti oleh peningkatan kebutuhan sumber
daya komputasi. Model dengan kemampuan yang lebih tinggi pada umumnya
membutuhkan kapasitas memori dan GPU yang lebih besar, sedangkan pada MAS
suatu persoalan dapat diselesaikan melalui beberapa kali proses inferensi.
Pada layanan komersial berbasis API, penggunaan model umumnya dikenai biaya
berdasarkan jumlah \textit{token} masukan dan keluaran \citep{openaipricing,
anthropicpricing}. Sementara itu, model \textit{open-weight} yang dijalankan
pada GPU lokal maupun sewaan menimbulkan biaya berdasarkan durasi
penggunaan sumber daya komputasi. Dengan demikian, pada kedua skema tersebut,
pengurangan jumlah \textit{token} yang diproses dan waktu inferensi berpotensi
menurunkan biaya penggunaan.

Pertimbangan lain dalam penggunaan LLM adalah privasi dan keamanan data,
karena penggunaan layanan tertutup mengharuskan data dikirim dan diproses
melalui infrastruktur pihak ketiga. Di antara organisasi yang mengalami
insiden keamanan terkait AI, 97\% menyatakan tidak memiliki kontrol akses AI
yang memadai, sementara 63\% belum memiliki kebijakan tata kelola AI
\citep{ibmbreach}. Model \textit{open-weight} yang dijalankan pada
infrastruktur yang dikendalikan sendiri menjadi salah satu alternatif yang
menarik, meskipun kemampuan model tersebut pada kondisi tertentu masih dapat
berada di bawah model komersial.

Pada MAS berbasis LLM, komunikasi antaragen umumnya berlangsung melalui teks,
yaitu keluaran satu agen dituliskan sebagai rangkaian \textit{token} diskret
sebelum digunakan sebagai masukan bagi agen berikutnya \citep{latentmas}.
Pendekatan tersebut memiliki keunggulan dari sisi keterbacaan dan kemudahan
interpretasi, tetapi memerlukan proses \textit{decoding} untuk menghasilkan
teks dan pemrosesan kembali ketika teks tersebut digunakan sebagai masukan
bagi agen berikutnya. Akibatnya, jumlah \textit{token} yang diproses dan waktu
inferensi dapat meningkat seiring dengan bertambahnya panjang penalaran dan
jumlah agen yang terlibat.

Sejumlah penelitian kemudian berupaya mengurangi ketergantungan terhadap
representasi token diskret pada dua bagian yang berbeda. Pada tingkat
penalaran dalam satu agen, \textit{Soft Thinking} \citep{softthinking} dan
LatentMAS \citep{latentmas} mengganti pemilihan token diskret dengan
representasi pada ruang kontinu atau ruang laten. Pada \textit{Soft Thinking},
masukan untuk langkah berikutnya dibentuk dari distribusi peluang token,
sedangkan pada LatentMAS representasi tersebut dapat dibentuk langsung dari
\textit{hidden state}. Pada tingkat
komunikasi antaragen, LatentMAS meneruskan KV-\textit{cache} sebagai
representasi keadaan internal dari satu agen ke agen berikutnya. Pendekatan-pendekatan tersebut berpotensi mengurangi
penggunaan \textit{token} dan waktu inferensi, tetapi menghilangkan atau
mengurangi representasi token diskret juga dapat memengaruhi akurasi sistem.

\textit{Mixture of Inputs} (MoI) \citep{moi} menempuh pendekatan yang berbeda
dengan mempertahankan kedua sumber informasi tersebut, yaitu distribusi
peluang keluaran model dan token yang terpilih melalui proses
\textit{sampling}. Dengan demikian, informasi mengenai kemungkinan token lain
yang tersedia dalam distribusi model tidak sepenuhnya hilang setelah satu
token dipilih. Penggabungan kedua sumber informasi tersebut diturunkan melalui
pendekatan inferensi Bayes.

Penelitian ini menguji kedua aspek tersebut secara bersamaan dalam satu
rancangan eksperimen. Pada aspek penalaran, MoI \citep{moi} ditetapkan sebagai
metode utama dan dibandingkan dengan tiga cara lain dalam membentuk
representasi langkah penalaran, yaitu \textit{raw} sebagaimana digunakan oleh
LatentMAS, \textit{soft}, dan \textit{sample} yang tetap memilih satu token.
Pada aspek komunikasi, komunikasi berbasis teks dibandingkan dengan komunikasi
berbasis KV-\textit{cache} sebagaimana digunakan oleh LatentMAS.

Evaluasi dilakukan menggunakan model \textit{open-weight} Qwen3-8B
\citep{qwen3} pada tiga \textit{benchmark}, yaitu GSM8K \citep{gsm8k},
ARC-Challenge \citep{arc}, dan HumanEval+ \citep{evalplus}. Seluruh perlakuan
dilakukan secara \textit{training-free} dan menggunakan konfigurasi bawaan
masing-masing metode tanpa pencarian parameter. Dengan demikian, penelitian
ini diarahkan untuk menjawab apakah Multi-Agent System dapat dibuat lebih
efisien dalam penggunaan \textit{token} dan waktu inferensi dengan tetap
mempertahankan tingkat akurasinya.


\section{Pembatasan Masalah}

Agar penelitian tetap terarah dan tidak meluas di luar ruang lingkup yang
ditetapkan, penelitian ini dibatasi pada hal-hal berikut.

\begin{enumerate}

\item Model bahasa yang digunakan adalah Qwen3-8B dengan inferensi lokal.
      Penelitian tidak menguji generalisasi terhadap ukuran maupun keluarga
      model bahasa lainnya.

\item Seluruh intervensi terhadap proses penalaran bersifat
      \textit{training-free}. Tidak dilakukan perubahan bobot model,
      \textit{fine-tuning}, maupun optimasi lainnya.

\item Metode penalaran laten yang dibandingkan dibatasi pada
      \textit{raw}, \textit{soft}, \textit{sample}, dan
      \textit{Mixture of Inputs} (MoI).

\item Medium komunikasi utama yang dibandingkan adalah komunikasi berbasis
      teks dan komunikasi laten berbasis KV-\textit{cache}.

\item Evaluasi kemampuan penalaran dibatasi pada GSM8K,
      ARC-Challenge, dan HumanEval-Plus. Ketiga tugas digunakan untuk
      merepresentasikan keluaran numerik, pilihan simbolik, dan
      penulisan kode.

\end{enumerate}


\section{Tujuan dan Manfaat Penelitian}

Sejalan dengan latar belakang yang telah diuraikan, penelitian ini bertujuan untuk:

\begin{enumerate}

      \item Merumuskan \textit{Mixture of Inputs} bersama tiga cara lain
            membentuk langkah penalaran laten, yaitu \textit{raw},
            \textit{soft}, dan \textit{sample}, ke dalam satu kerangka yang
            sama sehingga keempatnya dapat dibandingkan secara adil pada model
            Qwen3-8B, baik pada medium komunikasi berbasis teks maupun
            KV-\textit{cache}.

      \item Menganalisis efisiensi penggunaan token dan waktu inferensi dari
            masing-masing metode komunikasi dan penalaran laten, baik
            antara komunikasi berbasis teks dengan KV-\textit{cache} maupun antar-metode
            KV-\textit{cache} itu sendiri.

      \item Menganalisis hubungan antara efisiensi dan akurasi untuk
            mengetahui metode yang memberikan keseimbangan paling baik antara
            pengurangan biaya komputasi dan kemampuan penyelesaian tugas.

\end{enumerate}


Manfaat yang diharapkan adalah sebagai berikut.

\begin{enumerate}
    \item Memberikan kontribusi bagi kajian komunikasi antaragen berbasis LLM,
          khususnya mengenai penggunaan representasi kontinu dan ruang laten
          sebagai alternatif penalaran berbasis token diskret, sekaligus
          bahan pembanding bagi penelitian lanjutan yang mengembangkan metode
          komunikasi laten pada Multi-Agent System.
    \item Menyediakan perbandingan beberapa metode komunikasi dan representasi
          informasi dalam kondisi eksperimen yang terkendali, yaitu model dasar
          dan \textit{benchmark} yang sama sehingga karakteristik
          masing-masing metode dari segi akurasi, penggunaan token, dan waktu
          inferensi dapat dipahami.
    \item Mengetahui apakah pengurangan komunikasi berbasis teks melalui
          representasi laten, khususnya \textit{Mixture of Inputs}, dapat
          menghasilkan Multi-Agent System yang lebih cepat dan lebih murah
          tanpa mengorbankan akurasi secara signifikan, sebagai pertimbangan
          praktis dalam merancang sistem yang dijalankan pada model
          \textit{open-weight} di infrastruktur lokal maupun \textit{cloud}
          GPU.
\end{enumerate}

\section{Tinjauan Pustaka}

Berikut beberapa literatur kunci yang menjadi acuan penelitian ini.

\citet{latentmas} mengusulkan LatentMAS, kerangka kolaborasi
\textit{multi-agent} di ruang laten dengan dua mekanisme yang dipakai
penelitian ini. Pertama, menghasilkan pikiran laten di dalam satu agen: model
menambahkan representasi lapisan terakhir secara autoregresif sebagai pengganti
langkah pemilihan token, lalu menyelaraskannya kembali ke ruang \textit{embedding}
masukan melalui pemetaan linear berbasis regresi \textit{ridge}. Kedua,
pemindahan memori kerja antaragen: KV-\textit{cache} seluruh lapisan diteruskan
dari satu agen ke agen berikutnya tanpa penulisan teks. Penelitian ini memakai
mekanisme kedua sebagai medium komunikasi laten, dan mekanisme pertama sebagai
salah satu formulasi penalaran yang dibandingkan (\textit{raw}).

\citet{softthinking} mengusulkan \textit{soft thinking}, yang mengganti
pemilihan token diskret selama penalaran dengan \textit{concept token} berupa
kombinasi berbobot \textit{embedding} token menurut distribusi peluangnya.
\citet{stochsoft} menelaah mekanisme kerjanya dan menunjukkan bahwa penalaran
lunak yang deterministik cenderung terkunci pada token berpeluang tertinggi.
Penelitian ini membandingkan kedua sisi tersebut secara langsung: formulasi
\textit{soft} yang mempertahankan seluruh distribusi dan formulasi
\textit{sample} yang mengambil satu token melalui \textit{sampling}.

\citet{moi} mengusulkan \textit{Mixture of Inputs} dan memandang pembentukan
masukan bagi langkah berikutnya sebagai persoalan inferensi Bayes. Distribusi
peluang keluaran model diperlakukan sebagai parameter konsentrasi
\textit{prior} Dirichlet dan token hasil pengambilan sampel sebagai satu
observasi multinomial, sehingga melalui kekonjugatan Dirichlet--Multinomial
diperoleh \textit{posterior} berbentuk tertutup. Masukan bagi langkah
berikutnya dibentuk dari ekspektasi \textit{posterior} $\mathbb{E}[w\mid\cdot]$
atas bobot campuran tersebut, dengan besar kontribusi \textit{prior} relatif
terhadap observasi diatur oleh entropi distribusi, yaitu oleh tingkat
keyakinan model, sehingga metode tersebut menyesuaikan diri tanpa penalaan
parameter. MoI dirancang untuk penalaran satu model dan belum diuji sebagai
mekanisme komunikasi pada sistem \textit{multi-agent}; celah itulah yang
menjadi titik tolak penelitian ini.

\citet{coconut} dan survei \citet{latentsurvey} menempatkan ketiga pendekatan
di atas dalam paradigma penalaran laten yang lebih luas, yaitu pemindahan
langkah penalaran dari barisan token diskret ke ruang representasi kontinu.


\section{Metodologi Penelitian}

Metodologi yang digunakan adalah studi literatur dan eksperimen terkendali.
Studi literatur dilakukan dengan menelaah publikasi mengenai penalaran laten,
kolaborasi \textit{multi-agent}, dan evaluasi model bahasa, kemudian
menurunkan persamaan langkah laten tiap metode dari sumber aslinya ke dalam
satu kerangka notasi yang sama. Eksperimen dilakukan dengan menjalankan sistem
\textit{multi-agent} berbasis model \textit{open-weight} Qwen3-8B pada tolok
ukur GSM8K, ARC-Challenge, dan HumanEval+. Medium komunikasi antaragen dan
formulasi representasi laten diperlakukan sebagai faktor yang divariasikan,
sedangkan model, topologi agen, \textit{prompt}, banyaknya langkah laten,
suhu. Seluruh kondisi dievaluasi pada soal yang sama sehingga perbandingannya bersifat berpasangan.
% Analisis dilakukan menggunakan uji statistik nonparametrik yang sesuai untuk
% data berpasangan, disertai selang kepercayaan \textit{bootstrap} dan koreksi
% terhadap perbandingan berganda. Seluruh komputasi dilakukan menggunakan bahasa Python.

\section{Sistematika Penulisan}

Penyusunan laporan ini mengacu pada sistematika berikut.

\noindent\textsc{Bab I Pendahuluan}

Bab ini berisi latar belakang, pembatasan masalah, tujuan dan manfaat
penelitian, tinjauan pustaka, metodologi penelitian, dan sistematika
penulisan.

\noindent\textsc{Bab II Landasan Teori}

Bab ini menguraikan dasar teori yang digunakan, mencakup probabilitas dan
distribusi yang relevan, teori informasi, inferensi Bayes, aljabar linear dan
representasi vektor, arsitektur \textit{transformer} beserta KV-\textit{cache},
penalaran laten dan keempat formulasi pembentukan representasi laten, sistem
\textit{multi-agent}, konsep evaluasi empiris model bahasa, dan pengujian
statistik yang dipakai.

\noindent\textsc{Bab III Metodologi Penelitian}

Bab ini menetapkan objek penelitian berupa persamaan langkah laten,
menurunkan keempat persamaan yang dibandingkan ke dalam satu kerangka bersama,
menyatakan rancangan dua sumbu eksperimen, hipotesis, data dan tugas evaluasi,
protokol pengambilan sampel berpasangan, metrik, serta prosedur analisis
statistik.

\noindent\textsc{Bab IV Hasil dan Pembahasan}

Bab ini menyajikan hasil eksperimen beserta pembahasannya, mencakup
perbandingan akurasi antar-persamaan, keterbedaan hasil menurut jenis tugas,
biaya komputasi medium laten, fidelitas simbolik keluaran, bukti mekanisme
yang menjelaskan pola hasil, serta batas berlaku seluruh temuan.

\noindent\textsc{Bab V Kesimpulan dan Saran}

Bab ini memuat kesimpulan yang menjawab tujuan penelitian beserta saran bagi
penelitian selanjutnya.
